\documentclass[11pt]{article}
\usepackage[T1]{fontenc}
\usepackage{lmodern}
\usepackage[margin=1in]{geometry}
\usepackage{amsmath,amssymb,amsthm,mathtools}
\usepackage{microtype}
\usepackage[colorlinks=true,linkcolor=blue,citecolor=blue,urlcolor=blue]{hyperref}
\newtheorem{theorem}{Theorem}[section]
\newtheorem{lemma}[theorem]{Lemma}
\newtheorem{proposition}[theorem]{Proposition}
\newcommand{\R}{\mathbb R}
\newcommand{\dd}{\,dx}
\newcommand{\dt}{\,dt}
\newcommand{\grad}{\nabla}
\newcommand{\Div}{\operatorname{div}}
\newcommand{\supp}{\operatorname{supp}}
\DeclareMathOperator{\Lip}{Lip}
\title{An entropy chain rule on bounded Lipschitz domains}
\author{Kaiwen Jin\\\small\href{mailto:yui.ui.siki@gmail.com}{yui.ui.siki@gmail.com}}
\date{October 12, 2026}
\begin{document}
\maketitle
\begin{abstract}
We prove an exact entropy chain rule for nonnegative densities satisfying a
weak continuity equation with no boundary flux on any bounded Lipschitz
domain. Finite kinetic action and Fisher information, together with finite
initial entropy and weak continuity in time, imply absolute continuity of the
entropy on the closed time interval. The formula holds for the prescribed
flux, with vacuum allowed and without upper or lower density bounds. This
extends the convex-domain chain rule to nonsmooth nonconvex domains.
The proof reflects both the density and the flux across local boundary graphs.
A truncated entropy identity then permits the localization terms to cancel
before the logarithmic truncation is removed.
\end{abstract}

\section{The chain rule}\label{sec:theorem}
The entropy chain rule connects a weak continuity equation to the energy
balance of a diffusion equation. For example, in the Fokker--Planck equation
$\partial_t\rho=\Div(\grad\rho+\rho\grad V)$ the flux is
$F=-\grad\rho-\rho\grad V$. A chain rule for a general prescribed flux is
also needed when identifying limits that satisfy an energy-dissipation
inequality. Canc\`es, Monsaingeon and Natale prove such a rule on convex
domains in \cite[Proposition A.1]{CMN}; they ask immediately afterwards for
an extension to nonsmooth nonconvex Lipschitz domains.

Casteras, Flaim and Monsaingeon obtain an upper chain rule in smooth
nonconvex geometry \cite[Theorem 3]{CFM}. Their exact formula for the
optimal flux extends to a prescribed flux under an additional approximation
condition on the logarithmic gradient \cite[Theorem 4]{CFM}; the concluding
discussion of Section 4 of that version leaves Lipschitz geometry open.
A related chain rule on Lipschitz domains is proved by Erbar and Meglioli
\cite[Proposition 5.15]{EM} for Dirichlet boundary data. Its dissipation
functional incorporates the prescribed boundary trace
\cite[Definition 5.7]{EM}. Here the boundary condition is imposed only on
the flux, in the weak formulation below.

Let $d\geq1$, let $\Omega\subset\R^d$ be a bounded Lipschitz domain, and
let $T>0$. Put $Q=(0,T)\times\Omega$. For $V\in C^2(\overline\Omega)$
and $\pi=e^{-V}$, define
\begin{equation}\label{eq:entropy}
 E(\rho)=\int_\Omega\bigl(\beta(\rho)+\rho V\bigr)\dd,
 \qquad \beta(r)=r\log r-r+1,\qquad \beta(0)=1.
\end{equation}
The notation $L^1_w$ means $L^1$ with its weak topology. All gradients are
weak spatial gradients and all spacetime integrals are Lebesgue integrals.
We set $|F|^2/\rho=0$ where $(\rho,F)=(0,0)$ and $|F|^2/\rho=+\infty$
where $\rho=0$ and $F\ne0$.

\begin{theorem}\label{thm:main}
Suppose that $\rho\geq0$ belongs to $C([0,T];L^1_w(\Omega))$, with
$\int_\Omega\rho(t)\dd=M>0$ for every $t\in[0,T]$, and that
$F\in L^1(Q;\R^d)$. Assume the no-flux weak continuity equation
\begin{equation}\label{eq:ce}
 \int_0^T\!\int_\Omega
   \bigl(\rho\,\partial_t\varphi+F\cdot\grad\varphi\bigr)\dd\dt
 =\int_\Omega\rho(T)\varphi(T)\dd
  -\int_\Omega\rho(0)\varphi(0)\dd
\end{equation}
for every $\varphi\in C^1([0,T]\times\overline\Omega)$, and assume
$E(\rho(0))<\infty$ and
\begin{equation}\label{eq:energy}
 \int_0^T\!\int_\Omega\left(
   \frac{|F|^2}{\rho}
   +4\pi\left|\grad\sqrt{\rho/\pi}\right|^2\right)\dd\dt<\infty.
\end{equation}
The Fisher term in \eqref{eq:energy} requires
$\sqrt{\rho(t)/\pi}\in H^1(\Omega)$ for almost every $t$.
Then $E(\rho(t))$ is finite for every $t\in[0,T]$, is absolutely continuous
on $[0,T]$, and satisfies, for all $0\leq s\leq t\leq T$,
\begin{equation}\label{eq:main}
 E(\rho(t))-E(\rho(s))
 =\int_s^t\!\int_\Omega
   \frac{F}{\sqrt\rho}\cdot
       \bigl(2\grad\sqrt\rho+\sqrt\rho\,\grad V\bigr)\dd\dt.
\end{equation}
Here $F/\sqrt\rho=0$ on $\{\rho=0\}$. In particular,
\begin{equation}\label{eq:derivative}
 \frac{d}{dt}E(\rho(t))
 =\int_\Omega F(t)\cdot\grad\log(\rho(t)/\pi)\dd
 \quad\text{for almost every }t,
\end{equation}
where the integrand means the pairing in \eqref{eq:main}.
\end{theorem}

Thus the hypotheses and the prescribed-flux conclusion of
\cite[Proposition A.1]{CMN} hold on every bounded Lipschitz domain, including
those that are simultaneously nonsmooth and nonconvex. There is no density
bound or smooth-gradient approximation hypothesis.

Spatial convolution proves the formula away from the boundary. At a
Lipschitz graph, an even reflection of the density must be accompanied by a
corresponding reflection of the flux. The resulting change of variables
preserves the scalar entropy production even though it is not an orthogonal
reflection. We first keep the entropy derivative bounded, so the terms
created by a partition of unity are integrable and cancel. Only then do we
recover the logarithm. Sections \ref{sec:local}--\ref{sec:boundary} give these
two steps, and Section \ref{sec:limit} proves the theorem.

\section{A local truncated identity}\label{sec:local}
We use elementary Sobolev composition facts in the following forms: a
Sobolev gradient vanishes almost everywhere on a level set, and composition
with a bi-Lipschitz map preserves $W^{1,p}$ with the almost-everywhere chain
rule. These are \cite[Corollary 21 and Theorem 31]{H}; the absolute continuity
on lines characterization in \cite[Theorem 18]{H} also justifies the product
and reflection arguments below. Gradients here are column vectors, so for
$z\circ S$ the composition formula is
$\grad(z\circ S)=DS^T(\grad z)\circ S$.

Choose an even $\theta\in C^\infty(\R;[0,1])$ that equals one on
$[-1,1]$, vanishes outside $[-2,2]$, and is nonincreasing in $|a|$.
For $n\geq1$, set
\begin{equation}\label{eq:truncation}
 \kappa_n(r)=\theta(\log r/n)\quad(r>0),\qquad \kappa_n(0)=0,
 \qquad \beta_n''(r)=\frac{\kappa_n(r)}r,
 \qquad \beta_n(1)=\beta_n'(1)=0.
\end{equation}
We extend $\beta_n''$ by zero at $r=0$. It vanishes near zero and for large
$r$, so $\beta_n$ is smooth on $[0,\infty)$ and $\beta_n'$ is bounded for
fixed $n$. Also $\beta_n\geq0$.

\begin{lemma}\label{lem:local}
Let $U\subset\R^d$ be bounded and open. Suppose that
$r\geq0$ belongs to $C([0,T];L^1_w(U))$, with uniformly bounded mass,
$G\in L^1((0,T)\times U;\R^d)$, and
$\partial_t r+\Div G=0$ in distributions on $(0,T)\times U$.
Assume $\sqrt r\in L^2(0,T;H^1(U))$ and
\[
 \int_0^T\!\int_U\frac{|G|^2}{r}\dd\dt<\infty.
\]
Use the same vacuum conventions as in Theorem \ref{thm:main}.
For every compactly supported Lipschitz function $\psi$ on $U$ and every
$0\leq s\leq t\leq T$,
\begin{align}\label{eq:local}
 &\int_U\psi\beta_n(r(t))\dd-\int_U\psi\beta_n(r(s))\dd\notag\\
 &\qquad=\int_s^t\!\int_U\left[
    \psi\kappa_n(r)\frac{G}{\sqrt r}\cdot2\grad\sqrt r
       +\beta_n'(r)G\cdot\grad\psi\right]\dd\dt.
\end{align}
No boundary condition on $\partial U$ is required.
\end{lemma}

\begin{proof}
Put $v=\sqrt r$. The Sobolev product rule gives
$r\in L^1(0,T;W^{1,1}(U))$ and $\grad r=2v\grad v$.
The gradient $\grad v$ vanishes almost everywhere on $\{v=0\}$.
The action assumption also gives $G=0$ there. Hence
\[
 a=G/v,\qquad b=2\grad v
\]
are in $L^2((0,T)\times U)$, with $a=0$ on vacuum.

Let $\eta_\varepsilon$ be a nonnegative smooth spatial mollifier of unit
integral, supported in the $\varepsilon$-ball. Extend $r,G$ by zero solely
to define $r_\varepsilon=\eta_\varepsilon*r$ and
$G_\varepsilon=\eta_\varepsilon*G$.
On a compact set $K\Subset U$, and for sufficiently small $\varepsilon$,
\[
 \grad r_\varepsilon=\eta_\varepsilon*\grad r,
 \qquad \partial_t r_\varepsilon=-\Div G_\varepsilon.
\]
The kernels in these identities stay inside $U$.
Weighted Cauchy--Schwarz gives
\begin{equation}\label{eq:convolution-bounds}
 \frac{|G_\varepsilon|^2}{r_\varepsilon}
   \leq\eta_\varepsilon*|a|^2,
 \qquad
 \frac{|\grad r_\varepsilon|^2}{r_\varepsilon}
   \leq\eta_\varepsilon*|b|^2.
\end{equation}
If $r_\varepsilon=0$, both numerators vanish, and we set the normalized
vectors equal to zero.

At almost every point where $r>0$, spatial Lebesgue differentiation implies
that these normalized vectors converge to $a$ and $b$. At almost every
vacuum point, $|a|^2=|b|^2=0$, so \eqref{eq:convolution-bounds} proves their
convergence to zero as well. Spatial differentiation at almost every fixed
time, followed by Fubini, gives the spacetime assertion. The convolutions of
the zero-extended $|a|^2$ and $|b|^2$ converge in spacetime $L^1$ by
continuity of spatial translations. They are therefore uniformly integrable.
The bounds \eqref{eq:convolution-bounds} and almost-everywhere convergence
then imply
\begin{equation}\label{eq:normalized-convergence}
 \frac{G_\varepsilon}{\sqrt{r_\varepsilon}}\longrightarrow a,
 \qquad
 \frac{\grad r_\varepsilon}{\sqrt{r_\varepsilon}}\longrightarrow b
 \quad\text{strongly in }L^2((0,T)\times K).
\end{equation}
Indeed, each squared difference is uniformly integrable and tends to zero
almost everywhere, so its integral tends to zero.

For fixed $\varepsilon$, weak $L^1$ continuity of $r$ implies continuity of
$r_\varepsilon(t,x)$ in $t$ against each convolution kernel. Uniform mass
bounds give uniform spatial derivative bounds on $K$, hence continuity into
$C^k(K)$ for each fixed $k$. Moreover,
$\|\Div G_\varepsilon(t)\|_{C^k(K)}\leq C_{\varepsilon,k}\|G(t)\|_1$.
Integrating the distributional equation in time therefore gives an absolutely
continuous $C^k(K)$ representative, which agrees with the continuous one
at every time. Applying the ordinary chain rule and integrating by parts in
space yields
\begin{align}\label{eq:regularized}
 &\int_U\psi\beta_n(r_\varepsilon(t))\dd
       -\int_U\psi\beta_n(r_\varepsilon(s))\dd\notag\\
 &\quad=\int_s^t\!\int_U\left[
 \psi\beta_n''(r_\varepsilon)
                   G_\varepsilon\cdot\grad r_\varepsilon
   +\beta_n'(r_\varepsilon)G_\varepsilon\cdot\grad\psi
                                  \right]\dd\dt.
\end{align}
Take $K$ to contain the support of $\psi$ in its interior; integration by
parts is valid for this Lipschitz cutoff.

At each fixed endpoint, $r_\varepsilon\to r$ locally in $L^1$, and
$\beta_n$ is Lipschitz. The endpoint terms thus converge for every $s,t$.
In the first integrand of \eqref{eq:regularized}, write
\[
 \beta_n''(r_\varepsilon)G_\varepsilon\cdot\grad r_\varepsilon
 =\kappa_n(r_\varepsilon)
   \frac{G_\varepsilon}{\sqrt{r_\varepsilon}}\cdot
   \frac{\grad r_\varepsilon}{\sqrt{r_\varepsilon}}.
\]
The products converge in $L^1$ by \eqref{eq:normalized-convergence}.
The multiplier $\kappa_n(r_\varepsilon)$ is bounded and converges almost
everywhere, including on vacuum. For the second integrand,
$G_\varepsilon\to G$ in $L^1$, while $\beta_n'$ is bounded and continuous;
splitting off $G_\varepsilon-G$ and applying dominated convergence to the
remaining term proves $L^1$ convergence. This gives \eqref{eq:local}.
\end{proof}

\section{Reflection at the boundary}\label{sec:boundary}
We first extend the no-flux identity \eqref{eq:ce} from $C^1$ tests on the
closed domain to Lipschitz spatial tests. If $\zeta$ is
bounded and Lipschitz on $\R^d$, smooth mollifications converge uniformly
on $\overline\Omega$, their gradients are uniformly bounded, and those
gradients converge almost everywhere in $\Omega$. Dominated convergence
against $F\in L^1$ extends \eqref{eq:ce} to
$\varphi(t,x)=\alpha(t)\zeta(x)$, $\alpha\in C^1([0,T])$.
Consequently
\begin{equation}\label{eq:lip-tests}
 \int_\Omega\zeta\rho(t)\dd-\int_\Omega\zeta\rho(s)\dd
  =\int_s^t\!\int_\Omega F\cdot\grad\zeta\dd\dt
 \quad(0\leq s\leq t\leq T).
\end{equation}
The scalar curve on the left is continuous by weak $L^1$ continuity. Its
distributional derivative is the displayed $L^1$ function, so its absolutely
continuous representative agrees at every time. A bounded Lipschitz function
on $\overline\Omega$ can be extended to $\R^d$ and treated in the same way.

\begin{proposition}\label{prop:global}
Suppose that $\rho,F$ have the continuity and no-flux properties of
Theorem \ref{thm:main}, and that
$u=\sqrt\rho\in L^2(0,T;H^1(\Omega))$ and
$F/u\in L^2(Q;\R^d)$. Then, for every $n\geq1$ and every
$0\leq s\leq t\leq T$,
\begin{equation}\label{eq:global-truncated}
 \int_\Omega\beta_n(\rho(t))\dd-\int_\Omega\beta_n(\rho(s))\dd
 =\int_s^t\!\int_\Omega
        \kappa_n(\rho)\frac Fu\cdot2\grad u\dd\dt.
\end{equation}
\end{proposition}

\begin{proof}
Near a boundary point, take rigid coordinates in which
$\Omega$ lies above a Lipschitz graph $x_d=g(x')$.
By shrinking the chart, choose a ball $B\subset\R^{d-1}$ and $h>0$ so that
\begin{align*}
 U&=\{(x',x_d)\in B\times\R:\ |x_d-g(x')|<h\},\\
 U_+=U\cap\Omega&=\{(x',x_d)\in B\times\R:\ 0<x_d-g(x')<h\}
\end{align*}
lie within the graph neighborhood. Define
\begin{equation}\label{eq:reflection}
 R(x',x_d)=(x',2g(x')-x_d),\qquad
 A(x)=DR(x)=
 \begin{pmatrix} I_{d-1}&0\\2\grad g(x')^T&-1\end{pmatrix}
 \quad\text{almost everywhere}.
\end{equation}
The map $R$ is a bi-Lipschitz involution, $A(Rx)=A(x)$ and $A(x)^2=I$.
It preserves volume: this follows by the one-dimensional substitution in
$x_d$ for each $x'$ and Fubini. Its derivative is bounded in terms of
$\Lip(g)$. In dimension one, use ordinary reflection about an endpoint.

Extend the actual density and flux to $U$ by
\begin{equation}\label{eq:extension}
 r(x)=\rho(x),\quad G(x)=F(x)\quad(x\in U_+),\qquad
 r(Rx)=\rho(x),\quad G(Rx)=A(x)F(x)\quad(x\in U_+).
\end{equation}
The density extension is a bounded linear map on $L^1$, so it preserves weak
continuity in time and gives a uniform mass bound.
To check $\sqrt r\in L^2(0,T;H^1(U))$, flatten the graph by
$S(y',y_d)=(y',y_d+g(y'))$. Composition with $S$ maps $u|_{U_+}$ to
$H^1(B\times(0,h))$. Even reflection across $y_d=0$ belongs to
$H^1(B\times(-h,h))$: on almost every vertical line the one-dimensional
$H^1$ representative has a limit at zero and its even extension is
absolutely continuous; the tangential derivatives are reflected unchanged.
The absolute continuity on lines characterization then applies. Composing
back with $S^{-1}$ gives $\sqrt r$ and the gradient identity
\begin{equation}\label{eq:reflected-gradient}
 \grad\sqrt r(Rx)=A(x)^T\grad u(x)\quad\text{for almost every }x\in U_+.
\end{equation}
These composition facts also give the required time-integrated norm bound.
The reflected kinetic action is bounded because $G(Rx)/\sqrt{r(Rx)}=A(x)F(x)/u(x)$.
Only first derivatives of the graph occur.

For $\phi\in C_c^\infty(U)$, the folded spatial test
$\zeta(x)=\phi(x)+\phi(Rx)$ on $U_+$, extended by zero on the rest of
$\Omega$, is Lipschitz on $\overline\Omega$ and has an ambient Lipschitz
extension. Its support stays away from the outer chart boundary.
Using \eqref{eq:lip-tests} and changing variables on the lower half of $U$
proves $\partial_t r+\Div G=0$ in $(0,T)\times U$, since
\[
 \grad\zeta(x)=\grad\phi(x)+A(x)^T\grad\phi(Rx).
\]
Thus Lemma \ref{lem:local} applies to the reflected pair. No normal trace
of $F$ is needed.

Choose a finite smooth nonnegative ambient partition of unity $\chi_j$,
whose sum is one on a neighborhood of $\overline\Omega$.
Each support is compactly contained in an interior patch or a boundary
chart $U$ of the preceding kind. Such a partition follows from a finite
open cover of the compact closure by choosing smaller neighborhoods and
smooth bumps. In an interior patch apply \eqref{eq:local} with
$\psi=\chi_j$, $r=\rho$, $G=F$. For a boundary chart use the cutoff
\[
 \psi(x)=\chi_j(x)\quad(x\in U_+),\qquad
 \psi(Rx)=\chi_j(x)\quad(x\in U_+).
\]
It is compactly supported and Lipschitz on $U$, as follows by composition
with $(x',x_d)\mapsto(x',g(x')+|x_d-g(x')|)$, and
$\grad\psi(Rx)=A(x)^T\grad\chi_j(x)$.
Both pairings in \eqref{eq:local} are invariant under this reflection:
\begin{align}\label{eq:invariance}
 (AF/u)\cdot(2A^T\grad u)&=(F/u)\cdot2\grad u,\notag\\
 (AF)\cdot(A^T\grad\chi_j)&=F\cdot\grad\chi_j.
\end{align}
Indeed, $A^2=I$, although $A$ need not be orthogonal.
All three terms in \eqref{eq:local} are therefore twice their upper-half
integrals. Dividing by two gives the same formula as in an interior patch:
\begin{align*}
 &\int_\Omega\chi_j\beta_n(\rho(t))\dd
          -\int_\Omega\chi_j\beta_n(\rho(s))\dd\\
 &\quad=\int_s^t\!\int_\Omega\left[
   \chi_j\kappa_n(\rho)(F/u)\cdot2\grad u
       +\beta_n'(\rho)F\cdot\grad\chi_j\right]\dd\dt.
\end{align*}
Summing over $j$ cancels the last terms because
$\sum_j\grad\chi_j=0$ on $\Omega$. This proves
\eqref{eq:global-truncated}. The cancellation takes place for fixed $n$,
when $\beta_n'(\rho)F$ is integrable; it requires no integrability of
$\log(\rho)F$.
\end{proof}

\section{Removal of the truncation}\label{sec:limit}
\begin{proof}[Proof of Theorem \ref{thm:main}]
Recall $\pi=e^{-V}$, with $V\in C^2(\overline\Omega)$.
Set $w=\sqrt{\rho/\pi}$ and $u=\sqrt\rho=\sqrt\pi\,w$.
Since $\sqrt\pi$ and its reciprocal are bounded smooth multipliers,
\begin{equation}\label{eq:weight}
 2\grad u+u\grad V=2\sqrt\pi\,\grad w,
 \qquad
 4\int_Q|\grad u|^2\dd\dt
 \leq8\int_Q\pi|\grad w|^2\dd\dt
        +2TM\|\grad V\|_\infty^2.
\end{equation}
Also $\int_Q u^2\dd\dt=TM$. Thus $u\in L^2(0,T;H^1(\Omega))$,
and \eqref{eq:energy} gives $F/u\in L^2(Q)$, with $F=0$ on vacuum.
Proposition \ref{prop:global} applies.

The truncations increase to $\beta$. More explicitly,
\[
 \beta_n(r)=
 \begin{cases}
  \displaystyle\int_1^r(r-a)\frac{\kappa_n(a)}a\,da,&r\geq1,\\[6pt]
  \displaystyle\int_r^1(a-r)\frac{\kappa_n(a)}a\,da,&0\leq r\leq1.
 \end{cases}
\]
The kernels are nonnegative and $\kappa_n(a)\uparrow1$ for $a>0$.
Monotone convergence gives $\beta_n(r)\uparrow\beta(r)$, including
$\beta_n(0)\uparrow1$. The right-hand side of
\eqref{eq:global-truncated} is dominated by
$|F/u|\,2|\grad u|\in L^1(Q)$ and tends to the untruncated pairing.
Since $V$ is bounded and the mass is finite,
$\int_\Omega\beta(\rho(0))\dd<\infty$ follows from finite initial $E$.
Taking $s=0$ in \eqref{eq:global-truncated} first, monotone convergence
on the endpoints and dominated convergence on the right prove finiteness at
every time and
\[
 \int_\Omega\beta(\rho(t))\dd
  -\int_\Omega\beta(\rho(0))\dd
 =\int_0^t\!\int_\Omega(F/u)\cdot2\grad u\dd\dt.
\]
Subtracting two such identities gives the formula for every $s,t$, including
both endpoints of the actual weakly continuous curve.

Finally \eqref{eq:ce} with spatial test $V$ gives
\[
 \int_\Omega V\rho(t)\dd-\int_\Omega V\rho(s)\dd
 =\int_s^t\!\int_\Omega F\cdot\grad V\dd\dt.
\]
Adding the two formulas proves \eqref{eq:main}. By \eqref{eq:weight} its
integrand is
$(F/\sqrt\rho)\cdot2\sqrt\pi\,\grad\sqrt{\rho/\pi}$,
the product of the two $L^2$ factors in \eqref{eq:energy}.
It belongs to $L^1(Q)$. Hence \eqref{eq:main} proves absolute continuity
and differentiation gives \eqref{eq:derivative} almost everywhere.
\end{proof}

\paragraph{AI use.}
The proof development, source searches, drafting, and verification of this
note used OpenAI Codex. Automated reviews are not independent mathematical
certification. The manuscript is submitted for independent mathematical
review.

\begin{thebibliography}{9}
\bibitem{CMN}
C. Canc\`es, L. Monsaingeon and A. Natale,
\emph{Discretizing the Fokker--Planck equation with second-order accuracy:
a dissipation driven approach},
Numerische Mathematik \textbf{158} (2026), 1513--1563.
\href{https://doi.org/10.1007/s00211-026-01537-3}{doi:10.1007/s00211-026-01537-3}.
Preprint: \href{https://arxiv.org/abs/2410.03367v1}{arXiv:2410.03367v1}.
The cited proposition numbering refers to this preprint version.

\bibitem{CFM}
J.-B. Casteras, M. Flaim and L. Monsaingeon,
\emph{Entropy and Fisher information in non-convex domains: one chain to
rule them all},
\href{https://arxiv.org/abs/2512.05826v2}{arXiv:2512.05826v2},
revised September 16, 2026.

\bibitem{EM}
M. Erbar and G. Meglioli,
\emph{Gradient flow for a class of diffusion equations with Dirichlet
boundary data},
\href{https://arxiv.org/abs/2408.05987v2}{arXiv:2408.05987v2},
February 27, 2025.

\bibitem{H}
P. Haj\l{}asz,
\emph{Sobolev spaces, theory and applications},
lecture notes, Jyv\"askyl\"a Summer School, 1998.
\url{https://sites.pitt.edu/~hajlasz/Notatki/jyv-98c.pdf}.
\end{thebibliography}
\end{document}
